\section{결론}
표~\ref{tab:goal}에 목표 달성 현황을 정리하였다.

\begin{table*}[!t]
\centering\small
\caption{목표 달성 현황}\label{tab:goal}
\begin{tabularx}{\textwidth}{@{}p{40mm}p{28mm}Y@{}}
\toprule
목표 & 결과 & 근거 \\
\midrule
실패 발견율 30\% 이상 & \textbf{달성} & 로봇팔 58.3\%, 휴머노이드 40\% \\
Random 대비 20\% 이상 향상 & \textbf{달성} & 로봇팔 상대 250\% (16.7\%$\to$58.3\%), 휴머노이드 +40\%p \\
실패 유형 6종 중 4종 & \textbf{달성} & 충돌, 장애물 간섭, 목표 위치 점유, 작업자 안전 위험 + 도달 불가 후보 \\
3D 장면 구조화·후보 생성 & 구현 & 10축·21축 환경 변경, 장면 그래프·시뮬레이션 모델 생성, 실행 전 검사 \\
행동 기반 취약 조건 탐색 & 구현 & 관측$\to$가설$\to$제안$\to$FAIL 연결 확인 \\
회귀 테스트 자산화 & 구현 & 실행 기록 저장, 재실행·비교 실행기 \\
\bottomrule
\end{tabularx}
\end{table*}

본 과제는 3D 장면 파라미터화, LLM 기반 능동 실패 탐색, 단계적 판정, 회귀 자산화를 하나의
파이프라인으로 구현하였다. VLA 기반 로봇팔에서 AFS는 같은 예산으로 Random보다 3.5배 많은 실패를
찾았고, 휴머노이드에서도 Random이 찾지 못한 실패를 발견했다. 또한 장애물 높이와 목표 거리에 대한
PASS/FAIL 경계를 정량적으로 제시하여, 배포 전 검증에서 어떤 조건을 우선 점검해야 하는지 알려준다.

현재 실험은 시드 수와 반복 수가 작고, AFS가 동일 환경을 반복 제안하는 경향이 있다.
후속 과제로 (1) 더 많은 시드에서의 독립 평가, (2) 새로운 환경 탐색을 보상하는 다양성 목표,
(3) 유사 물체·카메라 가림 조건을 통한 인식 오류 탐색, (4) 선정 사례의 실제 회귀 재실행을 진행한다.
